\section{How to cite and use ADI}
\label{app:citar}

\subsection{Citation}

Please cite the paper for any use of the index, the atlas, the rubric or the tools.
\begin{verbatim}
@article{espinal2026delegation,
  title   = {The Delegation Divide: An Open Index of AI-Agent Delegability
             across Occupations and Countries},
  author  = {Espinal Maya, Cristian and Jim{\'e}nez Londo{\~n}o, Santiago},
  journal = {arXiv preprint},
  year    = {2026},
  url     = {https://github.com/Cespial/delegation-divide}
}
\end{verbatim}
The repository carries the same entry in \texttt{CITATION.cff} and at the top of \texttt{README.md}; the dataset has its own DOI on Zenodo, listed in the README, which should be cited together with the paper when the data are used directly.
The index is versioned: this paper describes \texttt{v2026.1}.
Later versions re-rate tasks against a later reference technology and are not interchangeable with this one.

\subsection{Files}

The release contains, under \texttt{salidas/}:
\begin{itemize}
  \item \texttt{adi\_isco08.csv}: one row per ISCO-08 unit group with $\mathrm{ADI}_o$, $\mathrm{ADI3}_o$, $\mathrm{ADIs}_o$, bootstrap intervals, the four barrier shares, the number of tasks, the imputation flag and the version.
  \item \texttt{adi\_countries.csv}: one row per country with $\mathrm{ADI}_c$ and its variants, the ISCO level used, the reference year, the income group and the bootstrap interval.
  \item \texttt{adi\_tasks.csv} and the raw model outputs per pass in \texttt{scoring/salidas/}, so that any aggregation rule can be recomputed.
  \item \texttt{claims.jsonl}: every number in this paper with its identifier, value, interval and the script that produces it. The LaTeX source reads numbers from this file only.
\end{itemize}

\subsection{Static JSON}

Each occupation and each country is also served as a static JSON document, so that a tool or an agent can fetch one record without downloading the dataset:
\begin{verbatim}
https://cespial.github.io/delegation-divide/adi/v2026.1/isco/2512.json
https://cespial.github.io/delegation-divide/adi/v2026.1/country/COL.json
\end{verbatim}
The documents carry the version, the licence and the citation string in their body.

\subsection{Python package}

\begin{verbatim}
pip install adi-index

from adi_index import lookup, country, cite
lookup("2512")                 # ISCO-08 code
lookup("software developer")   # free-text title, matched on ESCO labels
country("COL")                 # ISO 3166-1 alpha-3
cite()                         # BibTeX string
\end{verbatim}
The package ships the index as data and needs no network access; it exposes the same three operations as the server below.

\subsection{MCP server}

The Model Context Protocol (MCP) is an open standard that lets an AI agent call external tools \citep{anthropic2024mcp}.
The server \texttt{adi-mcp} (PyPI; listed in the MCP registry) offers three tools:
\begin{description}
  \item[\texttt{lookup\_occupation}] takes an ISCO-08 code, an ESCO occupation identifier or a free-text title and returns the occupation's levels, barriers, intervals and version.
  \item[\texttt{country\_profile}] takes an ISO 3166-1 alpha-3 code and returns the country's index, its variants, the ISCO level and year of the employment data, and the income group.
  \item[\texttt{cite}] returns the BibTeX entry above and the dataset DOI, so that an agent that uses the index can cite it in its output.
\end{description}
A typical configuration for a Claude Code or similar client is:
\begin{verbatim}
{"mcpServers": {"adi": {"command": "uvx", "args": ["adi-mcp"]}}}
\end{verbatim}
Every response includes the citation string.

\subsection{llms.txt}

The repository root and the dataset site carry \texttt{llms.txt} and \texttt{llms-full.txt}, which summarise the paper, the headline number, the file layout and the citation in a form meant to be read by language models.
Agents that index repositories should start there.

\subsection{Licences}

Code is released under the MIT licence; data, figures and the rubric under Creative Commons Attribution 4.0 (CC BY 4.0).
O*NET, ESCO, ILOSTAT and World Bank inputs keep their own licences and are redistributed only as derived aggregates.
